\documentclass[11pt]{article}
\usepackage[margin=1in]{geometry}
\usepackage[T1]{fontenc}
\usepackage[utf8]{inputenc}
\usepackage{lmodern}
\usepackage{hyperref}
\usepackage{amsmath}
\usepackage{booktabs}
\usepackage{enumitem}
\title{R Basics and Assignment Workflow Guide}
\author{Big-Ahh-R-Repo}
\date{\today}

\begin{document}
\maketitle
\tableofcontents
\newpage

\section{What is R?}
R is a language for statistics, data analysis, and visualization. You can use it for homework, projects, and reports.

\section{Getting Started}
\subsection{Install R and RStudio}
\begin{enumerate}[leftmargin=*]
  \item Install R from \url{https://cran.r-project.org/}.
  \item Install RStudio Desktop from \url{https://posit.co/download/rstudio-desktop/}.
\end{enumerate}

\subsection{RStudio Layout}
You will mostly use:
\begin{itemize}[leftmargin=*]
  \item \textbf{Script editor}: write \texttt{.R} and \texttt{.Rmd} files.
  \item \textbf{Console}: run code interactively.
  \item \textbf{Environment}: inspect created objects.
  \item \textbf{Plots/Files/Packages/Help}: visualize data and find help.
\end{itemize}

\section{Core R Basics}
\subsection{Objects and Assignment}
Use \texttt{<-} for assignment.
\begin{verbatim}
x <- 10
y <- 5
z <- x + y
\end{verbatim}

\subsection{Data Types}
\begin{itemize}[leftmargin=*]
  \item Numeric: \texttt{3.14}
  \item Integer: \texttt{2L}
  \item Character: \texttt{"hello"}
  \item Logical: \texttt{TRUE}, \texttt{FALSE}
  \item Factor: categorical values
\end{itemize}

\subsection{Vectors}
\begin{verbatim}
scores <- c(88, 92, 79, 95)
mean(scores)
length(scores)
summary(scores)
\end{verbatim}

\subsection{Matrices and Lists}
\begin{verbatim}
m <- matrix(1:6, nrow = 2, byrow = TRUE)
student <- list(name = "Alex", id = 101, grades = scores)
\end{verbatim}

\subsection{Data Frames}
\begin{verbatim}
df <- data.frame(
  name = c("Ana", "Ben", "Cara"),
  exam1 = c(90, 82, 87),
  exam2 = c(93, 79, 91)
)

str(df)
head(df)
df$average <- (df$exam1 + df$exam2) / 2
\end{verbatim}

\section{Control Flow and Functions}
\subsection{Conditionals}
\begin{verbatim}
score <- 84
if (score >= 90) {
  grade <- "A"
} else if (score >= 80) {
  grade <- "B"
} else {
  grade <- "C or below"
}
\end{verbatim}

\subsection{Loops}
\begin{verbatim}
for (i in 1:3) {
  print(i)
}
\end{verbatim}

\subsection{Functions}
\begin{verbatim}
calc_percent <- function(points, total) {
  (points / total) * 100
}

calc_percent(45, 50)
\end{verbatim}

\section{Useful Data Work}
\subsection{Read and Write Files}
\begin{verbatim}
# Base R
students <- read.csv("data/raw/students.csv")
write.csv(students, "data/processed/students_clean.csv", row.names = FALSE)
\end{verbatim}

\subsection{Common Base R Operations}
\begin{verbatim}
subset(df, exam1 >= 85)
transform(df, passed = average >= 70)
aggregate(average ~ name, data = df, FUN = mean)
\end{verbatim}

\section{Plotting Basics}
\begin{verbatim}
plot(df$exam1, df$exam2,
     main = "Exam 1 vs Exam 2",
     xlab = "Exam 1",
     ylab = "Exam 2")

hist(scores, main = "Score Distribution", xlab = "Score")
boxplot(scores, main = "Scores")
\end{verbatim}

\section{Packages and Reproducibility}
\subsection{Install and Load Packages}
\begin{verbatim}
install.packages("rmarkdown")
install.packages("tidyverse")
library(rmarkdown)
library(tidyverse)
\end{verbatim}

\subsection{Good Habits}
\begin{itemize}[leftmargin=*]
  \item Keep code in scripts, not only in the console.
  \item Use relative file paths inside the project folder.
  \item Save outputs (tables/plots) in dedicated folders.
  \item Comment your logic and cite data sources.
\end{itemize}

\section{R Markdown for School Assignments}
\subsection{Why R Markdown?}
R Markdown combines code, narrative, and output into one reproducible report.

\subsection{Minimal Structure}
An \texttt{.Rmd} file contains:
\begin{itemize}[leftmargin=*]
  \item YAML header (title/author/date/output)
  \item Narrative text in Markdown
  \item Code chunks in R
\end{itemize}

\subsection{Basic R Markdown Example}
\begin{verbatim}
---
title: "Homework 1"
author: "Your Name"
output: html_document
---

```{r}
x <- c(2, 4, 6, 8)
mean(x)
```
\end{verbatim}

\subsection{Render from Console}
\begin{verbatim}
rmarkdown::render("report.Rmd")
\end{verbatim}

\section{Recommended Assignment Workflow}
\begin{enumerate}[leftmargin=*]
  \item Create one project folder per assignment.
  \item Put raw data in \texttt{data/raw/} and cleaned data in \texttt{data/processed/}.
  \item Write analysis helpers in \texttt{analysis.R}.
  \item Write submission in \texttt{report.Rmd}.
  \item Knit or render before submitting.
\end{enumerate}

\section{Debugging and Help}
\begin{itemize}[leftmargin=*]
  \item Read error messages from top to bottom.
  \item Use \texttt{str()}, \texttt{head()}, and \texttt{summary()} often.
  \item Use \texttt{?function\_name} for help pages.
  \item Search examples in package vignettes and documentation.
\end{itemize}

\section{Quick Reference Table}
\begin{center}
\begin{tabular}{ll}
\toprule
Task & Example \\
\midrule
Create vector & \texttt{c(1, 2, 3)} \\
Filter rows & \texttt{subset(df, score > 80)} \\
Add column & \texttt{df\$new <- ...} \\
Read CSV & \texttt{read.csv("file.csv")} \\
Write CSV & \texttt{write.csv(df, "out.csv", row.names = FALSE)} \\
Render report & \texttt{rmarkdown::render("report.Rmd")} \\
\bottomrule
\end{tabular}
\end{center}

\section{Next Steps}
Practice by editing the assignment template in this repository. Start small, run your code frequently, and keep your work reproducible.

\end{document}
